% Moved verbatim from Paper A (new/sections/design.tex lines 27-90 (subjects with program rows) and 137-150, commit f96523f) on 2026-10-29. Raw material for Paper B; rewrite, do not reuse.
\subsection{Subjects}
Table~\ref{tab:subjects} lists the circuits and programs and the questions each serves. The rules are compared on
circuits, where an exact simulation gives the ground truth, and \tool{} is evaluated on programs, which it reads but
never runs. Corpus~A collects benchmark circuits of at most \numSizeAMaxQubits{} qubits from MQT
Bench~\cite{quetschlich2023mqt}, QASMBench~\cite{li2023qasmbench} and VeriQBench~\cite{chen2022veriqbench}. Every
qubit of every circuit is a criterion, which gives \numNPairs{} circuit--qubit pairs, \numNDynamicPairs{} of them in
dynamic circuits with mid-circuit measurement, reset or conditioned gates. The generated circuits come in families of known
structure that grow in depth and in width, corpus~L has wider circuits on which only some rules finish, and
QASMBench's medium and large sets have circuits far larger than corpus~A. The programs come from Bugs4Q, the Qiskit
textbook~\cite{qiskittextbook}, the corpus that LintQ's authors collected from GitHub and Qiskit HumanEval, and two
sets are our own, four parametric programs and ten programs with circuit loops.

\begin{table*}[t]
  \centering
  \caption{Circuits and programs of the study}
  \label{tab:subjects}
  \footnotesize\setlength{\tabcolsep}{4pt}\renewcommand{\arraystretch}{1.15}
  \newcommand{\sep}{\arrayrulecolor{black!25}\hline\arrayrulecolor{black}}
  \newcommand{\blk}[1]{\multicolumn{8}{@{}l}{\textit{#1}}\\}
  \begin{tabularx}{\textwidth}{@{}>{\raggedright\arraybackslash}X >{\raggedright\arraybackslash}p{3.05cm}
                  r r r r r >{\raggedright\arraybackslash}p{2.2cm}@{}}
    \toprule
    \textbf{Subject} & \textbf{Source} & \textbf{Count} & \textbf{Traced} & \textbf{Qubits} & \textbf{Operations} &
      \textbf{Code lines} & \textbf{Used for} \\
    \midrule
    \blk{Circuits}
    Corpus~A & MQT Bench~\cite{quetschlich2023mqt}, QASMBench~\cite{li2023qasmbench}, VeriQBench~\cite{chen2022veriqbench} &
      \numSubjACount{} & \numSubjATraced{} & \numSubjAMinQubits{}--\numSubjAMaxQubits{} & \numSubjAMedianOps{} / \numSubjAMaxOps{} & -- & RQ1 to RQ4 \\ \sep
    Generated circuits & ours & \numSubjGenCount{} & \numSubjGenTraced{} & \numSubjGenMinQubits{}--\numSubjGenMaxQubits{} &
      \numSubjGenMedianOps{} / \numSubjGenMaxOps{} & -- & RQ1, RQ4 \\ \sep
    Width sweeps & ours & \numSubjSweepCount{} & \numSubjSweepTraced{} & \numSubjSweepMinQubits{}--\numSubjSweepMaxQubits{} &
      \numSubjSweepMedianOps{} / \numSubjSweepMaxOps{} & -- & RQ4 \\ \sep
    Corpus~L & MQT Bench, QASMBench, VeriQBench & \numSubjLCount{} & \numSubjLTraced{} & \numSubjLMinQubits{}--\numSubjLMaxQubits{} &
      \numSubjLMedianOps{} / \numSubjLMaxOps{} & -- & RQ1, RQ4 \\ \sep
    QASMBench medium and large & QASMBench~\cite{li2023qasmbench} & \numSubjQbCount{} & \numSubjQbTraced{} &
      \numSubjQbMinQubits{}--\numSubjQbMaxQubits{} & \numSubjQbMedianOps{} / \numSubjQbMaxOps{} & -- & RQ4 \\
    \midrule
    \blk{Programs}
    Bugs4Q & Zhao et al.~\cite{zhao2023bugs4q} & \numSubjBugsCount{} & \numSubjBugsTraced{} &
      \numSubjBugsMinQubits{}--\numSubjBugsMaxQubits{} & \numSubjBugsMedianOps{} / \numSubjBugsMaxOps{} & \numSubjBugsMedianLines{} & RQ3, \tool{} \\ \sep
    Qiskit textbook & IBM~\cite{qiskittextbook} & \numSubjTextCount{} & \numSubjTextTraced{} &
      \numSubjTextMinQubits{}--\numSubjTextMaxQubits{} & \numSubjTextMedianOps{} / \numSubjTextMaxOps{} & \numSubjTextMedianLines{} & \tool{} \\ \sep
    LintQ's corpus & Paltenghi and Pradel~\cite{paltenghi2024lintq} & \numSubjLintqCount{} & \numSubjLintqTraced{} &
      \numSubjLintqMinQubits{}--\numSubjLintqMaxQubits{} & \numSubjLintqMedianOps{} / \numSubjLintqMaxOps{} & \numSubjLintqMedianLines{} & RQ4, LintQ comparison \\ \sep
    Qiskit HumanEval & Vishwakarma et al.~\cite{vishwakarma2024qhe} & \numSubjQheCount{} & \numSubjQheTraced{} &
      \numSubjQheMinQubits{}--\numSubjQheMaxQubits{} & \numSubjQheMedianOps{} / \numSubjQheMaxOps{} & \numSubjQheMedianLines{} & \tool{} \\ \sep
    Parametric programs, 5 sizes each & ours & \numSubjParamCount{} & \numSubjParamTraced{} of 20 &
      \numSubjParamMinQubits{}--\numSubjParamMaxQubits{} & \numSubjParamMedianOps{} / \numSubjParamMaxOps{} & \numSubjParamMedianLines{} & RQ4 \\ \sep
    Programs with circuit loops & ours & \numSubjCfCount{} & \numSubjCfTraced{} & \numSubjCfMinQubits{}--\numSubjCfMaxQubits{} &
      \numSubjCfMedianOps{} / \numSubjCfMaxOps{} & \numSubjCfMedianLines{} & Proposition~\ref{prop:loops} \\
    \bottomrule
  \end{tabularx}\\[3pt]
  \begin{minipage}{\textwidth}\footnotesize
  Traced counts the circuits read, of which QASMBench has \numSubjQbTraced{} in valid OpenQASM~2, and the programs from which
  \tool{} traces a circuit, or for the parametric programs the sizes, two of which loop beyond the interpreter's
  iteration limit. Qubits and operations describe the circuits, for programs the traced ones, as smallest to largest
  and as median / largest. Registers of unknown size are left out of the qubit range, and \numToolMaxOps{} operations is
  the interpreter's trace limit. Code lines exclude blank and comment lines and are given as medians. Faults are seeded
  into corpus~A, giving \numRqThreeUnits{} units from \numMutantsKilled{} mutants. Of Bugs4Q's bugs, \numBugsFourQIncluded{}
  port to circuits and \numBugsFourQKilled{} of them give \numLinesUnits{} units. A unit is a fault and a qubit whose
  measured distribution it changes.
  Programs by others are read and never run, and ours and our ports of the Bugs4Q bugs are also run.
  \end{minipage}
\end{table*}

\subsection{Cost, scalability and comparison with LintQ}
RQ4 times the rules on the generated family that grows in depth at a width of eight qubits, on width sweeps of two
generated families from \numSubjSweepMinQubits{} to \numSubjSweepMaxQubits{} qubits, and on corpus~L, each circuit in a fresh process with limits of \numCostWallLimit{}~s and 4~GiB. \tool{} is timed in
three settings, all exploratory. On LintQ's corpus it slices every qubit, up to 64, of each program that builds a
circuit, with a limit of 120~s per program. On QASMBench's medium and large circuits, read through the OpenQASM front
end, it slices every qubit in one pass, against three baselines that slice one qubit, Qiskit's DAG ancestors,
Qiskit's \texttt{LightCone} pass and the construction of our QPDG reimplementation, with a limit of 600~s per circuit.
On the parametric programs, a QFT, a GHZ preparation, a QAOA ring and a repetition code, a size constant grows the
trace while the source stays the same.

The comparison with LintQ uses LintQ's public data and never runs LintQ. For each of the \numLqQubitTp{} confirmed bugs
of its five analyses that concern qubits and bits, \tool{} traces the program and evaluates a predicate on its qubit
profiles, fixed before the comparison, at the warned line. The same predicates are evaluated on the warnings LintQ's
authors rejected and on all of its warnings of these analyses.
